%% toc: P3 31 31.5,31.6,31.7,31.8
\section{Proofs for two-cardinal models and cardinal transfer}
\label{sec:annex-p03-w03-p03-029}

\begin{annexitem}{annex:p03-w03-p03-029-theory-form}{thm:p03-w03-p03-029-theory-form}{From the relational source form to a countable theory}
\begin{proof}
Let \(\mathcal M\models T\) be a \((\kappa,\lambda)\)-model with
\(\kappa>\lambda\geq\aleph_0\).

\emph{Relational language.} Let \(\Lang'\) contain every relation symbol of
\(\Lang\), including \(U\), and, for each function symbol \(f\) of \(\Lang\)
of arity \(n\geq0\), a new relation symbol \(R_f\) of arity \(n+1\). Constant
symbols are the case \(n=0\). Since \(\Lang\) is countable, so is
\(\Lang'\). For an \(\Lang\)-structure \(\mathcal A\), let \(\mathcal A'\) be
the \(\Lang'\)-structure with the same universe, the same interpretation of
each relation symbol, and \(R_f^{\mathcal A'}\) the graph
\(\{(\bar a,b):f^{\mathcal A}(\bar a)=b\}\).

\emph{Term graphs.} For every \(\Lang\)-term \(t\) whose variables are among
\(\bar x\), and every variable \(y\) not occurring in \(t\), define an
\(\Lang'\)-formula \(\theta_t(\bar x,y)\) by recursion on \(t\):
\(\theta_{x_i}\) is \(x_i=y\), and \(\theta_{f t_1\ldots t_n}\) is
\[
  \exists z_1\cdots\exists z_n\Bigl(\theta_{t_1}(\bar x,z_1)\wedge\cdots
  \wedge\theta_{t_n}(\bar x,z_n)\wedge R_f z_1\ldots z_n y\Bigr)
\]
with \(z_1,\ldots,z_n\) new variables. For \(n=0\) this is \(R_f y\).

\emph{Claim 1.} For every \(\Lang\)-structure \(\mathcal A\), every
assignment \(\bar a\) to \(\bar x\) and every \(b\in\dom\mathcal A\),
\(\mathcal A'\models\theta_t[\bar a,b]\) if and only if
\(t^{\mathcal A}[\bar a]=b\). For a variable \(x_i\) both sides say
\(a_i=b\). For \(t=f t_1\ldots t_n\), by the induction hypothesis the only
witnesses for \(z_1,\ldots,z_n\) are \(c_k=t_k^{\mathcal A}[\bar a]\), and
\(R_f c_1\ldots c_n b\) holds in \(\mathcal A'\) exactly when
\(f^{\mathcal A}(c_1,\ldots,c_n)=b\), that is, when
\(t^{\mathcal A}[\bar a]=b\).

\emph{Formula translation.} Define \(\varphi^*\) by recursion. For an atomic
formula \(P t_1\ldots t_n\) with \(P\) a relation symbol, and for an
equation \(t_1=t_2\), put
\begin{align*}
  (P t_1\ldots t_n)^* &\defeq \exists z_1\cdots\exists z_n
     \Bigl(\textstyle\bigwedge_k\theta_{t_k}(\bar x,z_k)\wedge
     P z_1\ldots z_n\Bigr),\\
  (t_1=t_2)^* &\defeq \exists z_1\exists z_2
     \bigl(\theta_{t_1}(\bar x,z_1)\wedge\theta_{t_2}(\bar x,z_2)
     \wedge z_1=z_2\bigr).
\end{align*}
The translation commutes with \(\neg\), \(\wedge\) and \(\exists x\). The other connectives and \(\forall\) are treated through
their definitions from these.

\emph{Claim 2.} \(\mathcal A\models\varphi[\bar a]\) if and only if
\(\mathcal A'\models\varphi^*[\bar a]\). For atomic \(\varphi\), Claim~1
shows that the unique witnesses for the \(z_k\) are the values
\(t_k^{\mathcal A}[\bar a]\), so the translation holds exactly when the
relation or the equality holds of these values. The connective and
quantifier steps are immediate because \(\mathcal A\) and \(\mathcal A'\)
have the same universe.

\emph{Applying the source theorem.} \(\mathcal M'\) is a relational system
with a countable index set. \(U^{\mathcal M'}=U^{\mathcal M}\) is singulary,
and \(\aleph_0\leq\abs{U^{\mathcal M'}}=\lambda<\kappa=\abs{\dom\mathcal
M'}\). \Cref{thm:p03-w03-p03-029-mv-6-2} gives \(\mathcal B\equiv\mathcal
M'\) with \(\abs{\dom\mathcal B}=\aleph_1\) and
\(\abs{U^{\mathcal B}}=\aleph_0\). For each function symbol \(f\), the
sentence
\[
  \gamma_f\colon\quad
  \forall\bar x\,\exists y\bigl(R_f\bar x y\wedge
  \forall y'(R_f\bar x y'\to y'=y)\bigr)
\]
holds in \(\mathcal M'\), hence in \(\mathcal B\). So each
\(R_f^{\mathcal B}\) is the graph of a total function on
\(\dom\mathcal B\). Let \(\mathcal N\) be the \(\Lang\)-structure with
universe \(\dom\mathcal B\), the relations of \(\mathcal B\), and
\(f^{\mathcal N}(\bar a)\) the unique \(b\) with
\(R_f^{\mathcal B}(\bar a,b)\). Then \(\mathcal N'=\mathcal B\).

For every \(\Lang\)-sentence \(\varphi\), Claim~2 applied twice gives:
\(\mathcal M\models\varphi\) iff \(\mathcal M'\models\varphi^*\) iff
\(\mathcal B\models\varphi^*\) iff \(\mathcal N\models\varphi\). Hence
\(\mathcal N\equiv\mathcal M\), and in particular \(\mathcal N\models T\).
Finally \(\abs{\dom\mathcal N}=\aleph_1\) and
\(\abs{U^{\mathcal N}}=\abs{U^{\mathcal B}}=\aleph_0\), so \(\mathcal N\)
is an \((\aleph_1,\aleph_0)\)-model.
\end{proof}
\end{annexitem}

\begin{annexitem}{annex:p03-w03-p03-029-gap-one}{thm:p03-w03-p03-029-gap-one}{Gap-one reduction}
\begin{proof}
Put \(D=\dom\mathcal M\), \(Q=U^{\mathcal M}\) and \(\nu=\abs{D\setminus
Q}\). Since \(D=Q\cup(D\setminus Q)\) and \(\kappa\) is infinite,
\(\kappa=\max(\lambda,\nu)\). Because \(\lambda<\kappa\), this forces
\(\nu=\kappa\). As \(\kappa>\lambda\), we have \(\kappa\geq\lambda^+\), so
there is \(B\subseteq D\setminus Q\) with \(\abs B=\lambda^+\). Put
\(A=Q\cup B\). Then \(\abs A=\lambda+\lambda^+=\lambda^+\).

By the downward L\"owenheim--Skolem theorem (\cref{thm:downward-ls}), which
uses the Axiom of Choice, there is \(\mathcal N\preccurlyeq\mathcal M\) with
\(A\subseteq\dom\mathcal N\) and
\[
  \abs{\dom\mathcal N}\leq\max(\abs A,\abs\Lang,\aleph_0)
  =\max(\lambda^+,\aleph_0)=\lambda^+,
\]
using that \(\Lang\) is countable. Since \(A\subseteq\dom\mathcal N\),
\(\abs{\dom\mathcal N}\geq\lambda^+\), so equality holds. An elementary
substructure is a substructure, so \(U^{\mathcal N}=Q\cap\dom\mathcal N\),
and this equals \(Q\) because \(Q\subseteq A\). Thus \(\mathcal N\) is a
\((\lambda^+,\lambda)\)-model. Since \(\mathcal N\preccurlyeq\mathcal M\),
it is elementarily equivalent to \(\mathcal M\). For \(\lambda=\aleph_0\),
\(\lambda^+=\aleph_1\).
\end{proof}
\end{annexitem}

\begin{annexitem}{annex:p03-w03-p03-029-robinson}{thm:p03-w03-p03-029-robinson}{Robinson's example}
\begin{proof}
\emph{Item 1.} Let \(\varepsilon\) be the \(\Lang_E\)-sentence
\[
  \forall y\,\forall z\Bigl(\bigl(\neg Uy\wedge\neg Uz\wedge
  \forall x\,(Ux\to(Exy\leftrightarrow Exz))\bigr)\to y=z\Bigr).
\]
It holds in \(\mathcal R_X\). An element outside the predicate has the form
\((Y,1)\), and its \(E\)-predecessors in the predicate are the elements
\((x,0)\) with \(x\in Y\). Two such elements with the same predecessors
therefore have the same \(Y\). For each \(n\geq1\), the sentence stating
that at least \(n\) elements satisfy \(U\) holds in \(\mathcal R_X\) because
\(X\) is infinite.

Let \(\mathcal N\equiv\mathcal R_X\) and put \(V=U^{\mathcal N}\). The
sentences just listed hold in \(\mathcal N\), so \(V\) is infinite and
\(\mathcal N\models\varepsilon\). Define \(g\) on \(\dom\mathcal N\) by
\(g(a)=(a,0)\) for \(a\in V\), and
\(g(a)=(\{v\in V:\mathcal N\models E v a\},1)\) for \(a\notin V\). Then
\(g\) maps into \((V\times\{0\})\cup(\mathcal P(V)\times\{1\})\). It is
injective. Distinct elements of \(V\) have distinct images, and the tags
separate the two parts. Two elements outside \(V\) with the same image have
the same \(E\)-predecessors in \(V\), so they are equal by \(\varepsilon\).
Hence \(\abs{\dom\mathcal N}\leq\abs V+2^{\abs V}=2^{\abs V}\), the last
step because \(V\) is infinite and \(\abs V<2^{\abs V}\).

\emph{Item 2.} \(U^{\mathcal R_\omega}=\omega\times\{0\}\) has cardinality
\(\aleph_0\), and the universe has cardinality
\(\aleph_0+2^{\aleph_0}=2^{\aleph_0}>\aleph_0\).

\emph{Item 3.} A model of the theory of \(\mathcal R_X\) is elementarily
equivalent to \(\mathcal R_X\), because that theory is complete. If it
were a \((\kappa,\lambda)\)-model, item~1 would give
\(\kappa\leq2^\lambda\), contradicting \(\kappa>2^\lambda\).

\emph{Item 4.} Let \(\mathcal N\succcurlyeq\mathcal R_\omega\) with
\(\abs{\dom\mathcal N}>2^{\aleph_0}\). If \(U^{\mathcal N}\) were
countable, it would have cardinality \(\aleph_0\) by item~1, and item~1
would give \(\abs{\dom\mathcal N}\leq2^{\aleph_0}\). This is a
contradiction.

\emph{Item 5.} By item~2, \cref{thm:p03-w03-p03-029-gap-one} applies to
\(\mathcal R_\omega\) with \(\lambda=\aleph_0\), since the language is
finite. It gives an elementary substructure that is an
\((\aleph_1,\aleph_0)\)-model of the same theory.
\end{proof}
\end{annexitem}

\begin{annexitem}{annex:p03-w03-p03-029-hypotheses}{thm:p03-w03-p03-029-hypotheses}{Necessity of hypotheses}
\begin{proof}
\emph{Item 1.} For infinite \(\kappa\), the structure with universe \(\kappa\)
and \(U=\{0\}\) satisfies \(T_1\) and is a \((\kappa,1)\)-model. Every model
of \(T_1\) has a predicate of exactly one element, so none is an
\((\aleph_1,\aleph_0)\)-model.

\emph{Item 2.} For infinite \(\kappa\), the structure with universe \(\kappa\)
and \(U=\kappa\) satisfies \(T_2\) and is a \((\kappa,\kappa)\)-model. In
every model of \(T_2\) the predicate is the whole universe, so the two
coordinates coincide. In particular there is no
\((\aleph_1,\aleph_0)\)-model.

\emph{Item 3.} Take universe \(\omega_2\sqcup\omega\), with \(U\) the copy
of \(\omega\), and interpret \(c_\alpha\) as the element \(\alpha\) of the
copy of \(\omega_2\). This is a model of \(T_3\) of type
\((\aleph_2,\aleph_0)\). In every model of \(T_3\) the interpretations of
the \(c_\alpha\) are \(\aleph_2\) distinct elements, so the universe has
cardinality at least \(\aleph_2\). No model is an
\((\aleph_1,\aleph_0)\)-model.

In items~1 and~2 the language is countable; item~1 fails only the
hypothesis \(\lambda\geq\aleph_0\), item~2 only \(\kappa>\lambda\).
In item~3 the model is a two-cardinal model, and only the
countability of the language fails.
\end{proof}
\end{annexitem}

\begin{annexitem}{annex:p03-w03-p03-029-mv-3-1}{thm:p03-w03-p03-029-mv-3-1}{Finite subsystems of \(\mathfrak A^*\)}
\begin{proof}
Write \(\models_{\mathfrak A}\varphi[x_0,\ldots,x_{m-1}]\) for satisfaction
under the assignment of \(x_l\) to \(v_l\).

\emph{The equivalence.} The system \(\mathfrak A^*\mathbin{|}X\) has, for each
formula \(\varphi\) with least \(m\) as in
\cref{def:p03-w03-p03-029-mv-notions}, the relation \(\varphi^{\mathfrak
A}\cap X^m\). So \(f\) is an isomorphism of \(\mathfrak A^*\mathbin{|}X\) onto
\(\mathfrak B^*\mathbin{|}Y\) exactly when \(f\) is a bijection of \(X\) onto
\(Y\) and
\[
  (\dagger)\qquad
  \models_{\mathfrak A}\varphi[x_0,\ldots,x_{m-1}]
  \iff
  \models_{\mathfrak B}\varphi[f(x_0),\ldots,f(x_{m-1})]
\]
for every such \(\varphi\) and all \(x_0,\ldots,x_{m-1}\in X\). Since \(f\) is
onto \(Y\) by hypothesis, and \((\dagger)\) for \(v_0\approx v_1\) gives
injectivity, the isomorphism condition is equivalent to \((\dagger)\) alone.

A sentence of the language of \((\mathfrak A,x)_{x\in X}\) has the form
\(\psi(c_{x_0},\ldots,c_{x_{m-1}})\), where \(\psi\) is a formula of the
language of \(\mathfrak A\) with free variables among \(v_0,\ldots,v_{m-1}\)
and \(c_{x_l}\) is the constant for \(x_l\); renaming bound variables we may
assume that the substitution is free. It is true in
\((\mathfrak A,x)_{x\in X}\) iff \(\models_{\mathfrak
A}\psi[x_0,\ldots,x_{m-1}]\), and true in \((\mathfrak B,f(x))_{x\in X}\) iff
\(\models_{\mathfrak B}\psi[f(x_0),\ldots,f(x_{m-1})]\). If \(\psi\)'s least
\(m'\) is smaller than \(m\), the extra variables do not occur free and the
two conditions are the same as for \(m'\). Hence
\((\mathfrak A,x)_{x\in X}\equiv(\mathfrak B,f(x))_{x\in X}\) holds iff
\((\dagger)\) holds for every formula, that is, iff \(\mathfrak
A^*\mathbin{|}X\cong_f\mathfrak B^*\mathbin{|}Y\).

\emph{Item 1.} Every relation \(R_i^{\mathfrak A}\) of arity \(t\) is
\(\varphi^{\mathfrak A}\) for the atomic formula \(\varphi=P_i
v_0\cdots v_{t-1}\), so an automorphism of \(\mathfrak A^*\) is an
automorphism of \(\mathfrak A\). Conversely, an automorphism \(\sigma\) of
\(\mathfrak A\) satisfies \(\models_{\mathfrak A}\varphi[\bar x]\iff
\models_{\mathfrak A}\varphi[\sigma\bar x]\) for every formula: for atomic
formulas because \(\sigma\) preserves relations and equality, for
\(\neg\) and \(\wedge\) trivially, and for \(\exists v\) because \(\sigma\)
is onto, so a witness \(d\) on one side yields the witness \(\sigma(d)\),
and a witness \(e\) on the other side yields \(\sigma^{-1}(e)\). Hence
\(\sigma\) preserves every \(\varphi^{\mathfrak A}\). An elementary map is
injective: if \(x\neq x'\) in its domain, the sentence
\(\neg c_x\approx c_{x'}\) holds in \((\mathfrak A,x)_{x\in X}\), so
\(f(x)\neq f(x')\).

\emph{Item 2.} By the equivalence, the isomorphisms of finite subsystems of
\(\mathfrak A^*\) into \(\mathfrak A^*\) are exactly the elementary maps from
\(\mathfrak A\) to \(\mathfrak A\), and by item~1 their extensions to
automorphisms of \(\mathfrak A^*\) are their extensions to automorphisms of
\(\mathfrak A\).

\emph{Item 3.} \(\mathscr I\mathscr S_\omega(\mathfrak A^*)\subseteq\mathscr
I\mathscr S_\omega(\mathfrak B^*)\) says that for every finite non-empty
\(X\subseteq\abs{\mathfrak A}\) some finite \(Y\subseteq\abs{\mathfrak B}\)
and some \(f\) satisfy \(\mathfrak A^*\mathbin{|}X\cong_f\mathfrak
B^*\mathbin{|}Y\). By the equivalence, this is the existence of an elementary
map with domain \(X\).
\end{proof}
\end{annexitem}

\begin{annexitem}{annex:p03-w03-p03-029-mv-1-2}{thm:p03-w03-p03-029-mv-1-2}{Joint consistency}
\begin{proof}
\emph{(b).} Let \(\Lang\) be the language of \(\mathfrak D\) and
\(\mathfrak D'\), including constants for their distinguished elements. Add
new constants \(c_k\) for \(k\in K\) and new relation symbols \(Q_k\) for
\(k\in K'\), of the arity of \(S_k\); we may assume these new symbols are
pairwise distinct and distinct from those of \(\Lang\), as the source also
assumes \(K\cap K'=\emptyset\). Let
\(\Sigma=\mathrm{Th}(\mathfrak D,x_k)_{k\in K}\cup\mathrm{Th}(\mathfrak
D',S_k)_{k\in K'}\).

Let \(\Sigma_0\subseteq\Sigma\) be finite. The conjunction of the members of
\(\Sigma_0\cap\mathrm{Th}(\mathfrak D,x_k)_{k\in K}\) has the form
\(\varphi(c_{k_0},\ldots,c_{k_{n-1}})\) with \(\varphi\) an
\(\Lang\)-formula and \(k_0,\ldots,k_{n-1}\) distinct. Then
\(x_{k_0},\ldots,x_{k_{n-1}}\) satisfy \(\varphi\) in \(\mathfrak D\), so
\(\mathfrak D\) is a model of \(\exists v_0\cdots\exists
v_{n-1}\varphi\). This is an \(\Lang\)-sentence and \(\mathfrak
D\equiv\mathfrak D'\), so \(\mathfrak D'\) is a model of it; choose
witnesses \(e_0,\ldots,e_{n-1}\) in \(\mathfrak D'\). Expand
\((\mathfrak D',S_k)_{k\in K'}\) by interpreting \(c_{k_l}\) as \(e_l\) and the
other constants \(c_k\) arbitrarily. The expansion satisfies
\(\varphi(c_{k_0},\ldots,c_{k_{n-1}})\), and it satisfies every member of
\(\Sigma_0\cap\mathrm{Th}(\mathfrak D',S_k)_{k\in K'}\), since these
sentences contain no \(c_k\). So \(\Sigma_0\) has a model.

By the compactness theorem (\cref{thm:fol-compactness}) \(\Sigma\) has a
model \(\mathfrak E\). Let \(\mathfrak D''\) be its \(\Lang\)-reduct,
\(y_k=c_k^{\mathfrak E}\) and \(T_k=Q_k^{\mathfrak E}\). The theory of a
system contains each sentence of its language or its negation. Since
\((\mathfrak D'',y_k)_{k\in K}\) is a model of
\(\mathrm{Th}(\mathfrak D,x_k)_{k\in K}\), it satisfies no sentence outside
it, so \((\mathfrak D,x_k)_{k\in K}\equiv(\mathfrak D'',y_k)_{k\in K}\). In
the same way \((\mathfrak D',S_k)_{k\in K'}\equiv(\mathfrak
D'',T_k)_{k\in K'}\). Taking \(\Lang\)-sentences in the first equivalence
gives \(\mathfrak D''\equiv\mathfrak D\).

\emph{(a).} Apply (b) with \(K=\abs{\mathfrak D}\) and \(x_k=k\), and with
\(K'\) a copy of \(\abs{\mathfrak D'}\) disjoint from \(K\), indexed as
\(\{k':k\in\abs{\mathfrak D'}\}\), and \(S_{k'}=\{k\}\). Let \(\mathfrak D''\),
\(y_k\) and \(T_{k'}\) be as in (b). From
\((\mathfrak D,x)_{x\in\abs{\mathfrak D}}\equiv(\mathfrak
D'',y_x)_{x\in\abs{\mathfrak D}}\): the map \(x\mapsto y_x\) is injective
(sentences \(\neg c_x\approx c_{x'}\)), and for each relation \(R_i\) and
elements \(\bar x\) of \(\mathfrak D\), \(R_i^{\mathfrak D}\bar x\) holds iff
\(R_i^{\mathfrak D''}\bar y\) holds (atomic sentences and their negations);
distinguished elements are respected by the sentences \(d\approx c_x\). So
\(x\mapsto y_x\) is an isomorphism of \(\mathfrak D\) onto the subsystem of
\(\mathfrak D''\) on \(\{y_x\}\). From the second equivalence, each
\(T_{k'}\) is a singleton \(\{t_k\}\), by the sentence \(\exists
v\forall w(Q_{k'}w\leftrightarrow w\approx v)\). For a relation \(R_i\) of
arity \(t\) and elements \(k_1,\ldots,k_t\) of \(\mathfrak D'\), the sentence
\(\exists w_1\cdots\exists w_t(\bigwedge_lQ_{k_l'}w_l\wedge P_iw_1\cdots
w_t)\) holds in \((\mathfrak D',S)\) iff \(R_i^{\mathfrak D'}k_1\cdots k_t\),
and in \((\mathfrak D'',T)\) iff \(R_i^{\mathfrak D''}t_{k_1}\cdots
t_{k_t}\); the sentences \(\exists w\exists w'(Q_{k'}w\wedge
Q_{l'}w'\wedge\neg w\approx w')\) for \(k\neq l\) give injectivity, and
distinguished elements are treated in the same way. So \(k\mapsto t_k\) is an
isomorphism of \(\mathfrak D'\) onto the subsystem of \(\mathfrak D''\) on
\(\{t_k\}\). Hence \(\mathfrak D,\mathfrak D'\in\mathscr I\mathscr
S(\mathfrak D'')\).
\end{proof}
\end{annexitem}

\begin{annexitem}{annex:p03-w03-p03-029-mv-2-4d}{thm:p03-w03-p03-029-mv-2-4d}{One-point criterion for homogeneity}
\begin{proof}
\emph{Only if.} If \(\mathfrak A\) is homogeneous and \(f\) is an
elementary map with finite domain \(X\), then by
\cref{thm:p03-w03-p03-029-mv-3-1}(2) \(f\) extends to an automorphism
\(\sigma\) of \(\mathfrak A\). Put \(b=\sigma(a)\). By
\cref{thm:p03-w03-p03-029-mv-3-1}(1) \(\sigma\) preserves every formula, so
\((\mathfrak A,x)_{x\in X\cup\{a\}}\equiv(\mathfrak A,\sigma(x))_{x\in
X\cup\{a\}}\), and \(f\cup\{(a,b)\}=\sigma\mathbin{|}(X\cup\{a\})\) is
elementary.

\emph{If.} Assume the one-point condition. Two observations. First, if \(f\)
is an elementary map from \(\mathfrak A\) to \(\mathfrak A\), then \(f\) is
injective (\cref{thm:p03-w03-p03-029-mv-3-1}(1)) and \(f^{-1}\), on the
range of \(f\), is elementary, because \(\equiv\) is symmetric. Second,
the condition applied to \(f^{-1}\) and an element \(e\) gives \(a\) with
\(f^{-1}\cup\{(e,a)\}\) elementary, and then \(f\cup\{(a,e)\}\) is elementary
with \(e\) in its range.

Let \(f\) be an elementary map with finite non-empty domain, and enumerate
\(\abs{\mathfrak A}=\{e_n:n\in\omega\}\). Put \(f_0=f\). Given an elementary
\(f_{2n}\) with finite domain, let \(f_{2n+1}=f_{2n}\) if \(e_n\) is in its
domain, and otherwise \(f_{2n+1}=f_{2n}\cup\{(e_n,b)\}\) with \(b\) given by the
condition. Let \(f_{2n+2}=f_{2n+1}\) if \(e_n\) is in the range of
\(f_{2n+1}\), and otherwise \(f_{2n+2}=f_{2n+1}\cup\{(a,e_n)\}\) with \(a\)
given by the second observation. Each \(f_n\) is elementary with finite
domain, and \(f_n\subseteq f_{n+1}\). Let \(\sigma=\bigcup_n f_n\). Every
\(e_n\) lies in the domain of \(f_{2n+1}\) and in the range of
\(f_{2n+2}\), so \(\sigma\) is defined on all of \(\abs{\mathfrak A}\) and
onto it. It is injective, since any two points lie in the domain of one
injective \(f_n\). For a relation \(R_i\) of arity \(t\) and
\(x_0,\ldots,x_{t-1}\in\abs{\mathfrak A}\), choose \(n\) with all \(x_l\) in
the domain of \(f_n\); since \(f_n\) is elementary, \(R_i^{\mathfrak
A}x_0\cdots x_{t-1}\) holds iff \(R_i^{\mathfrak A}f_n(x_0)\cdots
f_n(x_{t-1})\) holds. So \(\sigma\) is an automorphism of \(\mathfrak A\)
extending \(f\), and \(\mathfrak A\) is homogeneous by
\cref{thm:p03-w03-p03-029-mv-3-1}(2).
\end{proof}
\end{annexitem}

\begin{annexitem}{annex:p03-w03-p03-029-mv-6-1}{thm:p03-w03-p03-029-mv-6-1}{Morley--Vaught Lemma 6.1}
\begin{proof}
Two facts are used repeatedly. (F1) If \(\mathfrak M\prec\mathfrak M'\)
then \(\mathfrak M\rceil J\prec\mathfrak M'\rceil J\), because every
sentence about the reducts with constants for elements of \(\mathfrak M\) is
a sentence about the full systems. (F2) If \(\mathfrak M\prec\mathfrak M'\)
and \(f\) is a function on a finite \(X\subseteq\abs{\mathfrak M}\) into
\(\abs{\mathfrak M}\), then \(f\) is an elementary map from \(\mathfrak M\)
to \(\mathfrak M\) iff it is one from \(\mathfrak M'\) to \(\mathfrak M'\),
since \((\mathfrak M,x)_{x\in X}\equiv(\mathfrak M',x)_{x\in X}\) and
\((\mathfrak M,f(x))_{x\in X}\equiv(\mathfrak M',f(x))_{x\in X}\).

\emph{Renaming.} If \(h\) is an injection of \(\abs{\mathfrak M}\) into
\(\abs{\mathfrak E}\) with \((\mathfrak M,m)_{m\in\abs{\mathfrak
M}}\equiv(\mathfrak E,h(m))_{m\in\abs{\mathfrak M}}\), there is
\(\mathfrak M'\succ\mathfrak M\) and an isomorphism \(h'\supseteq h\) of
\(\mathfrak M'\) onto \(\mathfrak E\). Indeed, choose a set \(D\) disjoint
from \(\abs{\mathfrak M}\) and a bijection \(e\) of \(D\) onto
\(\abs{\mathfrak E}\setminus h[\abs{\mathfrak M}]\); for instance
\(D=\{(u,\abs{\mathfrak M}):u\in\abs{\mathfrak E}\setminus
h[\abs{\mathfrak M}]\}\), which is disjoint from \(\abs{\mathfrak M}\) by
the axiom of foundation, since \((u,Z)\in Z\) would give the membership cycle
\(Z\in\{u,Z\}\in(u,Z)\in Z\). Let \(h'=h\cup e\) and give
\(\abs{\mathfrak M}\cup D\) the structure that makes \(h'\) an isomorphism
onto \(\mathfrak E\). The hypothesis on \(h\) says exactly that
\(\mathfrak M\prec\mathfrak M'\).

\emph{(a).} Let \(\mathfrak A_{\mathscr N}\) be an infinite system with
\(\mathscr N=\{\mathfrak B:\mathfrak B\equiv\mathfrak A_{\mathscr N}\}\).
The index set \(I_{\mathscr N}\) is countable, so the language is countable.

\emph{Step (one-point extension).} Let \(\mathfrak M\in\mathscr N\) be
countable, let \(p\) be an elementary map from \(\mathfrak M\rceil J\) to
\(\mathfrak M\rceil J\) with finite non-empty domain \(X\), and let
\(a\in\abs{\mathfrak M}\). We find a countable \(\mathfrak M'\succ\mathfrak
M\) and \(b\in\abs{\mathfrak M'}\) such that \(p\cup\{(a,b)\}\) is an
elementary map from \(\mathfrak M'\rceil J\) to \(\mathfrak M'\rceil J\).
Apply \cref{thm:p03-w03-p03-029-mv-1-2}(b) with
\(\mathfrak D=(\mathfrak M\rceil J,x)_{x\in X}\) and \(\mathfrak
D'=(\mathfrak M\rceil J,p(x))_{x\in X}\), the \(x\)-th distinguished element
of \(\mathfrak D'\) being \(p(x)\); \(\mathfrak D\equiv\mathfrak D'\) because
\(p\) is elementary. Take \(K=\{a\}\) with \(x_a=a\), and let the
relations \(S_k\) over \(\abs{\mathfrak D'}=\abs{\mathfrak M}\) be the
relations \(R_i^{\mathfrak M}\) for \(i\in I_{\mathscr N}\setminus J\) and
the singletons \(\{m\}\) for \(m\in\abs{\mathfrak M}\). Lemma~1.2(b) gives
\(\mathfrak D''=(\mathfrak G,z_x)_{x\in X}\), an element \(y\) and relations
\(T_i\) \((i\in I_{\mathscr N}\setminus J)\), \(T_{\{m\}}\) with
\begin{gather*}
  (\mathfrak M\rceil J,x,a)_{x\in X}\equiv(\mathfrak G,z_x,y)_{x\in X},\\
  (\mathfrak M\rceil J,p(x),R_i^{\mathfrak M},\{m\})\equiv
  (\mathfrak G,z_x,T_i,T_{\{m\}}).
\end{gather*}
By the second equivalence each \(T_{\{m\}}\) is a singleton \(\{t_m\}\),
and with \(\mathfrak E=(\mathfrak G,T_i)_{i\in I_{\mathscr N}\setminus J}\)
the map \(h(m)=t_m\) satisfies \((\mathfrak M,m)_{m}\equiv(\mathfrak
E,h(m))_{m}\), because a sentence \(\psi(c_{m_1},\ldots,c_{m_r})\) holds in
\((\mathfrak M,m)_m\) iff \(\exists\bar w(\bigwedge_lQ_{\{m_l\}}w_l\wedge
\psi(\bar w))\) holds in \((\mathfrak D',S_k)\), with each \(P_i\),
\(i\notin J\), read as the symbol for \(R_i^{\mathfrak M}\); moreover \(z_x=t_{p(x)}\), because the sentence saying that
the \(x\)-th distinguished element lies in the relation for \(\{p(x)\}\) holds
on the left. By renaming, transporting \(z_x\) and \(y\) along the isomorphism, we may
take \(\mathfrak M\prec\mathfrak E\) with \(h\) the identity, so
\(z_x=p(x)\). The first equivalence then reads
\((\mathfrak M\rceil J,x,a)_{x\in X}\equiv(\mathfrak E\rceil
J,p(x),y)_{x\in X}\), and by (F1) the left side is elementarily equivalent to
\((\mathfrak E\rceil J,x,a)_{x\in X}\). So \(p\cup\{(a,y)\}\) is an elementary
map from \(\mathfrak E\rceil J\) to \(\mathfrak E\rceil J\). By the downward
L\"owenheim--Skolem theorem (\cref{thm:downward-ls}) there is a countable
\(\mathfrak M'\prec\mathfrak E\) containing \(\abs{\mathfrak M}\cup\{y\}\);
then \(\mathfrak M\prec\mathfrak M'\) (both are elementary in
\(\mathfrak E\)), and by (F1), (F2) the map \(p\cup\{(a,y)\}\) is elementary
from \(\mathfrak M'\rceil J\) to \(\mathfrak M'\rceil J\). Put \(b=y\). If
\(a\in X\), the first equivalence contains \(c_a\approx c_x\) for \(x=a\), so
\(y=p(a)\) and \(p\cup\{(a,b)\}=p\).

This step is the instance, at power \(\omega\), of condition (3) in the
source's proof of Theorem~2.6(a) together with hypothesis (ii) of
Theorem~2.8\('\): the extension stays in \(\mathscr N\), while the map
concerns only the \(J\)-reduct.

\emph{Recursion (Theorem~2.6(a) at cofinality \(\omega\)).} Let
\(\pi(i,k)=\tfrac12(i+k)(i+k+1)+k\), a bijection of \(\omega\times\omega\)
onto \(\omega\) with \(\pi(i,k)\geq i\). By the downward
L\"owenheim--Skolem theorem choose a countable \(\mathfrak M_0\prec\mathfrak
A_{\mathscr N}\); it is infinite, since \(\mathfrak A_{\mathscr N}\) satisfies
the sentences saying that there are at least \(n\) elements. Suppose
\(\mathfrak M_0\prec\cdots\prec\mathfrak M_n\) are countable members of
\(\mathscr N\). When \(\mathfrak M_n\) is defined, fix an enumeration
\((p_{n,k},a_{n,k})_{k\in\omega}\), with repetitions allowed, of all pairs
\((p,a)\) with \(p\) an elementary map from \(\mathfrak M_n\rceil J\) to
\(\mathfrak M_n\rceil J\) with finite non-empty domain and
\(a\in\abs{\mathfrak M_n}\). There are countably many such pairs and at least
one (the identity on \(\{a\}\)). Let \(\pi(i,k)=n\); then \(i\leq n\), so
\(\mathfrak M_i\prec\mathfrak M_n\) and by (F1), (F2) the map \(p_{i,k}\) is
elementary from \(\mathfrak M_n\rceil J\) to \(\mathfrak M_n\rceil J\). Let
\(\mathfrak M_{n+1}\) and \(b\) be given by the step for \(\mathfrak M_n\),
\(p_{i,k}\), \(a_{i,k}\).

Let \(\mathfrak M=\bigcup_n\mathfrak M_n\). By the elementary chain theorem
(\cref{thm:elementary-chain}) \(\mathfrak M_n\prec\mathfrak M\) for every
\(n\), so \(\mathfrak M\in\mathscr N\); \(\mathfrak M\) is countable and
infinite, so denumerable. Let \(\mathfrak A=\mathfrak M\rceil J\in\mathscr
N\rceil J\). We check the condition of
\cref{thm:p03-w03-p03-029-mv-2-4d}. Let \(p\) be an elementary map from
\(\mathfrak A\) to \(\mathfrak A\) with finite non-empty domain \(X\), and
\(a\in\abs{\mathfrak A}\). The finite set \(X\cup p[X]\cup\{a\}\) lies in
some \(\abs{\mathfrak M_i}\), and by (F1), (F2) \(p\) is elementary from
\(\mathfrak M_i\rceil J\) to \(\mathfrak M_i\rceil J\), so
\((p,a)=(p_{i,k},a_{i,k})\) for some \(k\). At stage \(n=\pi(i,k)\) an
element \(b\) was found with \(p\cup\{(a,b)\}\) elementary from
\(\mathfrak M_{n+1}\rceil J\) to itself, and by (F1), (F2) it is elementary
from \(\mathfrak A\) to \(\mathfrak A\). Hence \(\mathfrak A\)
is homogeneous.

\emph{(b) (the argument of Theorem~2.3(b)).} Let \(\mathfrak A\),
\(\mathfrak A'\) be as in (b). By \cref{thm:p03-w03-p03-029-mv-3-1}(3),
each finite non-empty \(X\subseteq\abs{\mathfrak A}\) is the domain of an
elementary map from \(\mathfrak A\) to \(\mathfrak A'\), and symmetrically.

\emph{Forth.} Let \(f\) be an elementary map from \(\mathfrak A\) to
\(\mathfrak A'\) with finite non-empty domain \(X\), and
\(a\in\abs{\mathfrak A}\). Choose an elementary map \(g\) from
\(\mathfrak A\) to \(\mathfrak A'\) with domain \(X\cup\{a\}\). Then
\((\mathfrak A',g(x))_{x\in X}\equiv(\mathfrak A,x)_{x\in X}\equiv(\mathfrak
A',f(x))_{x\in X}\), and \(g\) is injective, so \(q=f\circ g^{-1}\) on
\(g[X]\) is an elementary map from \(\mathfrak A'\) to \(\mathfrak A'\). As
\(\mathfrak A'\) is homogeneous, \(q\) extends to an automorphism \(\sigma\)
of \(\mathfrak A'\) (\cref{thm:p03-w03-p03-029-mv-3-1}(2)). Put
\(b=\sigma(g(a))\). Then \(f\cup\{(a,b)\}=\sigma\circ g\), and since
\(\sigma\) preserves every formula,
\((\mathfrak A,x)_{x\in X\cup\{a\}}\equiv(\mathfrak A',g(x))_{x\in
X\cup\{a\}}\equiv(\mathfrak A',\sigma g(x))_{x\in X\cup\{a\}}\). So
\(f\cup\{(a,b)\}\) is elementary. \emph{Back.} The same argument with the roles
of \(\mathfrak A\) and \(\mathfrak A'\) exchanged, applied to \(f^{-1}\) and
using the homogeneity of \(\mathfrak A\), gives for each \(e\in\abs{\mathfrak
A'}\) an \(a\) with \(f\cup\{(a,e)\}\) elementary.

Enumerate \(\abs{\mathfrak A}=\{a_n\}\) and \(\abs{\mathfrak A'}=\{a'_n\}\).
Start with an elementary map \(f_0\) with domain \(\{a_0\}\), and alternate
forth steps for \(a_n\) and back steps for \(a'_n\) as in the proof of
\cref{thm:p03-w03-p03-029-mv-2-4d}. The union of the resulting chain is a
bijection of \(\abs{\mathfrak A}\) onto \(\abs{\mathfrak A'}\) that
preserves every relation, that is, an isomorphism. The hypothesis
\(\mathfrak A\equiv\mathfrak A'\) is implied by the equality of the classes
\(\mathscr I\mathscr S_\omega\) and is not used separately.

\emph{(c) (by Corollary~2.4(d)).} Let \(\mathfrak A=\bigcup_n\mathfrak A_n\).
It is countable as a countable union of countable sets and infinite, so
denumerable. By the elementary chain theorem \(\mathfrak A_n\prec\mathfrak
A\) for every \(n\). Let \(p\) be an elementary map from \(\mathfrak A\) to
\(\mathfrak A\) with finite non-empty domain \(X\), and
\(a\in\abs{\mathfrak A}\). Choose \(n\) with \(X\cup p[X]\cup\{a\}\subseteq
\abs{\mathfrak A_n}\). By (F2) \(p\) is elementary from \(\mathfrak A_n\) to
\(\mathfrak A_n\); it extends to an automorphism \(\sigma\) of the
homogeneous \(\mathfrak A_n\). With \(b=\sigma(a)\), \(p\cup\{(a,b)\}\) is
elementary from \(\mathfrak A_n\) to \(\mathfrak A_n\), hence, by (F2), from
\(\mathfrak A\) to \(\mathfrak A\). By \cref{thm:p03-w03-p03-029-mv-2-4d},
\(\mathfrak A\) is homogeneous.
\end{proof}
\end{annexitem}

\begin{annexitem}{annex:p03-w03-p03-029-mv-7}{thm:p03-w03-p03-029-mv-7}{Display (7), countable index set}
\begin{proof}
Let \(\Lang\) be the (countable) language of \(\mathfrak A\), with relation
symbols \(P_i\) for \(i\in I_{\mathfrak A}\), and let \(\Lang_J\) be the
language of the relativized reducts, with a relation symbol \(Q_j\) for each
\(j\in J\). For a system \(\mathfrak M\) similar to \(\mathfrak A\) with
\(R_{j_k}^{\mathfrak M}\neq\emptyset\), write \(\mathfrak E_k(\mathfrak M)
=\langle R_{j_k}^{\mathfrak M},R_{f_k(j)}^{\mathfrak M}\rangle_{j\in J}\) for
\(k=1,2\).

\emph{Relativization.} For an \(\Lang_J\)-formula \(\psi\) define the
\(\Lang\)-formula \(\psi^{(k)}\): replace each \(Q_j\) by \(P_{f_k(j)}\),
leave equations unchanged, commute with \(\neg\) and \(\wedge\), and put
\((\exists u\,\psi)^{(k)}=\exists u\,(P_{j_k}u\wedge\psi^{(k)})\). By
induction on \(\psi\): for every such \(\mathfrak M\) and elements
\(\bar a\) of \(R_{j_k}^{\mathfrak M}\),
\[
  (\ddagger)\qquad
  \mathfrak E_k(\mathfrak M)\models\psi[\bar a]
  \iff
  \mathfrak M\models\psi^{(k)}[\bar a].
\]
For atomic \(\psi\) this holds because the \(j\)-th relation of
\(\mathfrak E_k(\mathfrak M)\) is \(R_{f_k(j)}^{\mathfrak M}\) restricted to
\(R_{j_k}^{\mathfrak M}\) and \(\bar a\) lies in \(R_{j_k}^{\mathfrak M}\); the
connective steps are immediate; for \(\exists u\), the witnesses on the left
are the elements of \(R_{j_k}^{\mathfrak M}\), which are the witnesses
allowed by \(P_{j_k}u\) on the right. Consequences: (R1) if
\(\mathfrak M\prec\mathfrak M'\) then \(\mathfrak E_k(\mathfrak
M)\subseteq\mathfrak E_k(\mathfrak M')\) and \(\mathfrak E_k(\mathfrak
M)\prec\mathfrak E_k(\mathfrak M')\), by \((\ddagger)\) on both sides;
(R2) the hypothesis \(\mathfrak E_1(\mathfrak A)\equiv\mathfrak
E_2(\mathfrak A)\) says that \(\theta^{(1)}\leftrightarrow\theta^{(2)}\) holds
in \(\mathfrak A\) for every \(\Lang_J\)-sentence \(\theta\), so
\(\mathfrak E_1(\mathfrak M)\equiv\mathfrak E_2(\mathfrak M)\) for every
\(\mathfrak M\equiv\mathfrak A\).

\emph{Finite case.} Suppose \(R_{j_1}^{\mathfrak A}=\{a_1,\ldots,a_n\}\)
with the \(a_l\) distinct. The \(\Lang_J\)-sentence saying that there are
exactly \(n\) elements holds in \(\mathfrak E_1(\mathfrak A)\), hence in
\(\mathfrak E_2(\mathfrak A)\). Take \(\mathfrak B=\mathfrak A\). For finite
\(J_0\subseteq J\) let \(B(J_0)\) be the set of bijections of
\(R_{j_1}^{\mathfrak A}\) onto \(R_{j_2}^{\mathfrak A}\) that are isomorphisms
of the \(J_0\)-reducts. \(B(J_0)\neq\emptyset\): let \(\delta(w_1,\ldots,w_n)\)
be the conjunction of \(\neg w_l\approx w_{l'}\) for \(l\neq l'\) and, for each
\(j\in J_0\) and each tuple of indices of the arity of \(Q_j\), the atomic
formula \(Q_j w_{l_1}\cdots\) or its negation according as it holds of the
corresponding \(a\)'s in \(\mathfrak E_1(\mathfrak A)\). This is a finite
conjunction. \(\exists\bar w\,\delta\) holds in \(\mathfrak E_1(\mathfrak
A)\), hence in \(\mathfrak E_2(\mathfrak A)\); witnesses \(b_1,\ldots,b_n\) are
distinct, so they exhaust the \(n\)-element set \(R_{j_2}^{\mathfrak A}\),
and \(a_l\mapsto b_l\) lies in \(B(J_0)\). If \(J_0\subseteq J_1\) then
\(B(J_1)\subseteq B(J_0)\). All \(B(J_0)\) are subsets of the finite set of
bijections, so choose a finite \(J_0\) with \(\abs{B(J_0)}\) least. For any
finite \(J_1\), \(\emptyset\neq B(J_0\cup J_1)\subseteq B(J_0)\), so by
minimality \(B(J_0\cup J_1)=B(J_0)\) and \(B(J_0)\subseteq B(J_1)\). An element
of \(B(J_0)\) therefore preserves every \(Q_j\), and is an isomorphism of
\(\mathfrak E_1(\mathfrak A)\) onto \(\mathfrak E_2(\mathfrak A)\).

\emph{Infinite case: realization.} Let \(\mathfrak M\) be countable and
\(\Sigma(w_1,\ldots,w_r)\) a set of \(\Lang\)-formulas with parameters from
\(\abs{\mathfrak M}\) every finite subset of which is satisfied in
\(\mathfrak M\) by some tuple. Then some countable \(\mathfrak
M'\succ\mathfrak M\) contains a tuple satisfying all of \(\Sigma\). Indeed,
let \(\Gamma\) be the theory of \((\mathfrak M,m)_{m\in\abs{\mathfrak M}}\)
together with \(\Sigma(d_1,\ldots,d_r)\) for new constants \(d_l\). A finite
subset of \(\Gamma\) holds in \((\mathfrak M,m)_m\) expanded by a tuple
satisfying the finitely many members of \(\Sigma\) it contains. By
compactness (\cref{thm:fol-compactness}; the language is countable) and the
downward L\"owenheim--Skolem theorem, \(\Gamma\) has a countable model
\(\mathfrak E\); the map \(m\mapsto c_m^{\mathfrak E}\) satisfies the
hypothesis of the renaming in the proof of \cref{thm:p03-w03-p03-029-mv-6-1},
which yields \(\mathfrak M'\succ\mathfrak M\) isomorphic to the
\(\Lang\)-reduct of \(\mathfrak E\) by a map fixing \(\abs{\mathfrak M}\); the
preimages of the \(d_l^{\mathfrak E}\) satisfy \(\Sigma\).

\emph{Infinite case: tasks.} Suppose \(R_{j_1}^{\mathfrak A}\) is infinite.
Then \(\mathfrak A\) is infinite, and so is \(R_{j_2}^{\mathfrak A}\), by
(R2) and the sentences saying that there are at least \(n\) elements. These
facts hold in every \(\mathfrak M\equiv\mathfrak A\). For countable
\(\mathfrak M\equiv\mathfrak A\), a \emph{task} of \(\mathfrak M\) is one of:
\begin{enumerate}
  \item[(T1)] a triple \((k,p,a)\) with \(k\in\{1,2\}\), \(p\) an elementary
        map from \(\mathfrak E_k(\mathfrak M)\) to \(\mathfrak
        E_k(\mathfrak M)\) with finite non-empty domain, and
        \(a\in R_{j_k}^{\mathfrak M}\);
  \item[(T2)] a pair \((k,\bar a)\) with \(k\in\{1,2\}\) and \(\bar a\) a
        finite non-empty tuple from \(R_{j_k}^{\mathfrak M}\).
\end{enumerate}
A countable \(\mathfrak M\) has countably many tasks. \emph{Handling} a task
in a countable \(\mathfrak M'\succ\mathfrak M\) means: for (T1), finding
\(b\in R_{j_k}^{\mathfrak M'}\) with \(p\cup\{(a,b)\}\) elementary from
\(\mathfrak E_k(\mathfrak M')\) to itself; for (T2), with \(k'\neq k\),
finding \(\bar b\) from \(R_{j_{k'}}^{\mathfrak M'}\) with
\((\mathfrak E_k(\mathfrak M'),\bar a)\equiv(\mathfrak E_{k'}(\mathfrak
M'),\bar b)\), the tuples indexing the same constants. Every task of
\(\mathfrak M\) is a task of each \(\mathfrak M''\succ\mathfrak M\), by
(R1) and the argument of (F2) in the proof of
\cref{thm:p03-w03-p03-029-mv-6-1}.

A task \((k,p,a)\) of a countable \(\mathfrak M\) is handled in some
countable \(\mathfrak M'\succ\mathfrak M\). Let \(X\) be the domain of \(p\)
and \(\bar x\) an enumeration of \(X\). Let \(\Sigma(w)\) consist of
\(P_{j_k}w\) and all \(\psi^{(k)}(w,p\bar x)\) for \(\Lang_J\)-formulas
\(\psi(w,\bar u)\) with \(\mathfrak E_k(\mathfrak M)\models\psi[a,\bar x]\).
A finite subset is implied by \(P_{j_k}w\wedge\psi^{(k)}(w,p\bar x)\) for
one such \(\psi\) (a conjunction). \(\mathfrak E_k(\mathfrak
M)\models\exists w\,\psi[\bar x]\), so, \(p\) being elementary,
\(\mathfrak E_k(\mathfrak M)\models\exists w\,\psi[p\bar x]\), and by
\((\ddagger)\) a witness in \(R_{j_k}^{\mathfrak M}\) satisfies the finite
subset in \(\mathfrak M\). Realization gives \(\mathfrak M'\) and \(b\). For
every \(\psi\): if \(\mathfrak E_k(\mathfrak M)\models\psi[a,\bar x]\) then
\(\mathfrak E_k(\mathfrak M')\models\psi[b,p\bar x]\) by \((\ddagger)\); if
not, the same holds for \(\neg\psi\). Since \(\mathfrak E_k(\mathfrak
M)\prec\mathfrak E_k(\mathfrak M')\) by (R1), the truth of \(\psi[a,\bar
x]\) is the same in both, so \(p\cup\{(a,b)\}\) is elementary from
\(\mathfrak E_k(\mathfrak M')\) to itself.

A task \((k,\bar a)\) of a countable \(\mathfrak M\equiv\mathfrak A\) is
handled in some countable \(\mathfrak M'\succ\mathfrak M\). Let
\(\Sigma(\bar w)\) consist of the \(P_{j_{k'}}w_l\) and all
\(\psi^{(k')}(\bar w)\) with \(\mathfrak E_k(\mathfrak
M)\models\psi[\bar a]\). For one such \(\psi\),
\(\mathfrak E_k(\mathfrak M)\models\exists\bar w\,\psi\), so by (R2)
\(\mathfrak E_{k'}(\mathfrak M)\models\exists\bar w\,\psi\), and by
\((\ddagger)\) a witness from \(R_{j_{k'}}^{\mathfrak M}\) satisfies the
corresponding finite subset of \(\Sigma\) in \(\mathfrak M\). Realization
gives \(\mathfrak M'\) and \(\bar b\) with \((\mathfrak E_{k'}(\mathfrak
M'),\bar b)\equiv(\mathfrak E_k(\mathfrak M),\bar a)\equiv(\mathfrak
E_k(\mathfrak M'),\bar a)\), the last by (R1).

\emph{Infinite case: the chain.} With \(\pi\) as in the proof of
\cref{thm:p03-w03-p03-029-mv-6-1}, choose a countable \(\mathfrak
M_0\prec\mathfrak A\), and when \(\mathfrak M_n\) is defined fix an
enumeration \((\tau_{n,k})_{k\in\omega}\) of its tasks. If \(\pi(i,k)=n\), let
\(\mathfrak M_{n+1}\succ\mathfrak M_n\) be countable and handle
\(\tau_{i,k}\), a task of \(\mathfrak M_n\) since \(i\leq n\). Let
\(\mathfrak B=\bigcup_n\mathfrak M_n\). By the elementary chain theorem
\(\mathfrak M_n\prec\mathfrak B\), so \(\mathfrak B\equiv\mathfrak A\), and
\(\mathfrak B\) is countable. For \(k=1,2\), \(\mathfrak E_k(\mathfrak B)\)
is infinite and countable, hence denumerable, and \(\mathfrak
E_k(\mathfrak M_n)\prec\mathfrak E_k(\mathfrak B)\) by (R1).

\(\mathfrak E_k(\mathfrak B)\) is homogeneous: given \(p\) elementary from
\(\mathfrak E_k(\mathfrak B)\) to itself with finite non-empty domain and
\(a\in R_{j_k}^{\mathfrak B}\), all elements involved lie in some
\(\mathfrak M_i\), so \((k,p,a)\) is a task \(\tau_{i,l}\) of \(\mathfrak
M_i\). It is handled in \(\mathfrak M_{n+1}\), \(n=\pi(i,l)\), and the
resulting \(p\cup\{(a,b)\}\) is elementary from \(\mathfrak
E_k(\mathfrak B)\) to itself because \(\mathfrak E_k(\mathfrak
M_{n+1})\prec\mathfrak E_k(\mathfrak B)\).
\Cref{thm:p03-w03-p03-029-mv-2-4d} applies.

\(\mathscr I\mathscr S_\omega(\mathfrak E_1(\mathfrak B)^*)=\mathscr I\mathscr
S_\omega(\mathfrak E_2(\mathfrak B)^*)\): a finite non-empty
\(X\subseteq R_{j_k}^{\mathfrak B}\), enumerated as \(\bar a\), lies in
some \(\mathfrak M_i\); the task \((k,\bar a)\) is handled in some
\(\mathfrak M_{n+1}\) by \(\bar b\), and as above the equivalence
\((\mathfrak E_k(\mathfrak B),\bar a)\equiv(\mathfrak E_{k'}(\mathfrak
B),\bar b)\) holds; distinct entries of \(\bar a\) have distinct partners, so
this defines an elementary map from \(\mathfrak E_k(\mathfrak B)\) to
\(\mathfrak E_{k'}(\mathfrak B)\) with domain \(X\).
\Cref{thm:p03-w03-p03-029-mv-3-1}(3) gives both inclusions. By (R2)
\(\mathfrak E_1(\mathfrak B)\equiv\mathfrak E_2(\mathfrak B)\). Lemma~6.1(b)
(\cref{thm:p03-w03-p03-029-mv-6-1}) gives \(\mathfrak E_1(\mathfrak
B)\cong\mathfrak E_2(\mathfrak B)\).
\end{proof}
\end{annexitem}

\begin{annexitem}{annex:p03-w03-p03-029-mv-6-2}{thm:p03-w03-p03-029-mv-6-2}{Morley--Vaught Theorem 6.2}
\begin{proof}
The proof follows the source \cite[p.~55, proof of
Theorem~6.2]{MorleyVaught1962}, with each step written out. Let
\(\Lang\) be the language of \(\mathfrak A\), with relation symbols \(P_i\)
for \(i\in I_{\mathfrak A}\), and put \(\mu=\overline{R_j^{\mathfrak A}}\), so
\(\omega\leq\mu<\overline{\mathfrak A}\).

\emph{Step 1.} By the downward L\"owenheim--Skolem theorem
(\cref{thm:downward-ls}) with \(A=R_j^{\mathfrak A}\) there is
\(\mathfrak B_1\prec\mathfrak A\) with \(R_j^{\mathfrak
A}\subseteq\abs{\mathfrak B_1}\) and \(\overline{\mathfrak
B_1}\leq\max(\mu,\aleph_0)=\mu\); also \(\overline{\mathfrak B_1}\geq\mu\).
So \(\overline{\mathfrak B_1}=\mu<\overline{\mathfrak A}\), and
\(\abs{\mathfrak B_1}\neq\abs{\mathfrak A}\).

\emph{Step 2: display (7).} Let \(\mathfrak A^+=(\mathfrak
A,\abs{\mathfrak A},\abs{\mathfrak B_1})\), with new singulary indices
\(u\) (for \(\abs{\mathfrak A}\)) and \(c\) (for \(\abs{\mathfrak B_1}\)); its
index set is countable. Apply \cref{thm:p03-w03-p03-029-mv-7} with
\(j_1=u\), \(j_2=c\), \(J=I_{\mathfrak A}\) and \(f_1=f_2\) the inclusion
of \(I_{\mathfrak A}\) into \(I_{\mathfrak A}\cup\{u,c\}\). The relativized
reducts are \(\langle\abs{\mathfrak A},R_i^{\mathfrak A}\rangle=\mathfrak A\)
and \(\langle\abs{\mathfrak B_1},R_i^{\mathfrak A}\rangle=\mathfrak
A\mathbin{|}\abs{\mathfrak B_1}=\mathfrak B_1\), and \(\mathfrak
A\equiv\mathfrak B_1\) since \(\mathfrak B_1\prec\mathfrak A\). So there
is \(\mathfrak B^+\equiv\mathfrak A^+\) whose two relativized reducts are
isomorphic. The sentence \(\forall x\,P_ux\) holds in \(\mathfrak A^+\), so
\(\mathfrak B^+=(\mathfrak B_2,\abs{\mathfrak B_2},C)\) for a system
\(\mathfrak B_2\) and a set \(C\); thus \((\mathfrak B_2,C)\equiv(\mathfrak
A,\abs{\mathfrak B_1})\) and \(\mathfrak B_2\cong_f\mathfrak B_2\mathbin{|}C\)
for some \(f\).

\emph{Step 3: Lemma 6.1(a).} Let \(\mathscr N\) be the elementary type of
\((\mathfrak B_2,C,f)\), with the index \(c\) for \(C\) and a new index
\(\varrho\) for the graph of \(f\). It is infinite, as \(\mathfrak B_2\equiv\mathfrak
A\), and its index set is countable. By
\cref{thm:p03-w03-p03-029-mv-6-1}(a) with \(J=I_{\mathfrak A}\) there is
\((\mathfrak A_1,A_0,G)\in\mathscr N\) with \(\mathfrak A_1\) denumerable
and homogeneous. The following sentences hold in \((\mathfrak B_2,C,f)\):
those of \((\mathfrak B_2,C)\), which are those of \((\mathfrak
A,\abs{\mathfrak B_1})\), and those expressing that \(f\) is an isomorphism
onto \(\mathfrak B_2\mathbin{|}C\). So they hold in \((\mathfrak
A_1,A_0,G)\):
\begin{enumerate}
  \item[(s1)] \(G\) is the graph of an injective function \(g\) on
        \(\abs{\mathfrak A_1}\) onto \(A_0\), and for each \(i\in
        I_{\mathfrak A}\) of arity \(t\),
        \(\forall\bar x\forall\bar y(\bigwedge_{l<t}P_\varrho x_ly_l\to(P_i\bar
        x\leftrightarrow P_i\bar y))\). Hence \(\mathfrak A_1\cong_g\mathfrak
        A_0\), where \(\mathfrak A_0=\mathfrak A_1\mathbin{|}A_0\).
  \item[(s2)] \(\exists x\,\neg P_cx\), true since \(\abs{\mathfrak
        B_1}\neq\abs{\mathfrak A}\); so \(A_0\neq\abs{\mathfrak A_1}\).
  \item[(s3)] \(\forall x(P_jx\to P_cx)\); so \(R_j^{\mathfrak
        A_0}=R_j^{\mathfrak A_1}\cap A_0=R_j^{\mathfrak A_1}\).
  \item[(s4)] For each \(\Lang\)-formula \(\varphi(v_0,\ldots,v_{m-1})\),
        \(\forall\bar v(\bigwedge_lP_cv_l\to(\varphi\leftrightarrow\varphi^{(c)}))\),
        where \(\varphi^{(c)}\) is the relativization to \(P_c\) of the proof
        of \cref{thm:p03-w03-p03-029-mv-7} (with \(J=I_{\mathfrak A}\) and
        the identity for \(f_k\)). It holds in \((\mathfrak A,\abs{\mathfrak
        B_1})\) by \((\ddagger)\) there, because \(\mathfrak
        B_1\prec\mathfrak A\). In \((\mathfrak A_1,A_0)\), again by
        \((\ddagger)\), it says \(\mathfrak A_0\prec\mathfrak A_1\).
  \item[(s5)] For each \(n\), \(P_j\) has at least \(n\) elements; so
        \(R_j^{\mathfrak A_1}\) is infinite, and countable as a subset of
        \(\abs{\mathfrak A_1}\).
\end{enumerate}
So \(\mathfrak A_0\) and \(\mathfrak A_1\) are isomorphic, \(\mathfrak A_1\)
is a proper elementary extension of \(\mathfrak A_0\), and
\(R_j^{\mathfrak A_0}=R_j^{\mathfrak A_1}\) is denumerable. Moreover
\(\mathfrak A_1\equiv\mathfrak B_2\equiv\mathfrak A\), and \(\mathfrak A_0\),
being isomorphic to \(\mathfrak A_1\), is denumerable and homogeneous.

\emph{Step 4: the recursion.} We define \(\mathfrak A_\xi\) for
\(\xi<\omega_1\), extending the given \(\mathfrak A_0,\mathfrak A_1\), with
\[
  \mathfrak A_\xi\cong\mathfrak A_0,\qquad
  R_j^{\mathfrak A_\xi}=R_j^{\mathfrak A_0},\qquad
  \mathfrak A_\eta\prec\mathfrak A_\xi\text{ and }\mathfrak A_\eta\neq\mathfrak
  A_\xi\text{ for }\eta<\xi .
\]
These hold for \(\xi\leq1\) by Step~3.

\emph{Successor.} Given \(\mathfrak A_\xi\) with \(\xi\geq1\), choose an
isomorphism \(h\) of \(\mathfrak A_0\) onto \(\mathfrak A_\xi\). Choose a set
\(D\) disjoint from \(\abs{\mathfrak A_\xi}\) and a bijection \(e\) of
\(\abs{\mathfrak A_1}\setminus A_0\) onto \(D\) (as in the renaming in the
proof of \cref{thm:p03-w03-p03-029-mv-6-1}). Then \(h'=h\cup e\) is a
bijection of \(\abs{\mathfrak A_1}\) onto \(\abs{\mathfrak A_\xi}\cup D\); let
\(\mathfrak A_{\xi+1}\) be the system on \(\abs{\mathfrak A_\xi}\cup D\) that
makes \(h'\) an isomorphism of \(\mathfrak A_1\) onto it. This is the
system related to \(\mathfrak A_\xi\) as \(\mathfrak A_1\) is to
\(\mathfrak A_0\). Since \(\mathfrak A_0=\mathfrak A_1\mathbin{|}A_0\) and
\(h'\mathbin{|}A_0=h\), we get \(\mathfrak A_{\xi+1}\mathbin{|}\abs{\mathfrak
A_\xi}=\mathfrak A_\xi\). For a formula \(\varphi\) and \(\bar a\) from
\(\mathfrak A_\xi\): \(\mathfrak A_\xi\models\varphi[\bar a]\) iff
\(\mathfrak A_0\models\varphi[h^{-1}\bar a]\) iff \(\mathfrak
A_1\models\varphi[h^{-1}\bar a]\) iff \(\mathfrak
A_{\xi+1}\models\varphi[\bar a]\), using \(\mathfrak A_0\prec\mathfrak A_1\)
and \(h'h^{-1}\bar a=\bar a\). So \(\mathfrak A_\xi\prec\mathfrak A_{\xi+1}\),
and the extension is proper because \(D\neq\emptyset\) by (s2). Further
\(R_j^{\mathfrak A_{\xi+1}}=h'[R_j^{\mathfrak A_1}]=h[R_j^{\mathfrak
A_0}]=R_j^{\mathfrak A_\xi}=R_j^{\mathfrak A_0}\) by (s3), and
\(\mathfrak A_{\xi+1}\cong\mathfrak A_1\cong\mathfrak A_0\). For
\(\eta<\xi\), \(\mathfrak A_\eta\prec\mathfrak A_\xi\prec\mathfrak
A_{\xi+1}\), and \(\abs{\mathfrak A_\eta}\subseteq\abs{\mathfrak
A_\xi}\neq\abs{\mathfrak A_{\xi+1}}\).

\emph{Limit.} Let \(\delta<\omega_1\) be a limit and
\(\mathfrak A_\delta=\bigcup(\mathfrak A_\eta/\eta<\delta)\). By the
elementary chain theorem \(\mathfrak A_\eta\prec\mathfrak A_\delta\) for
\(\eta<\delta\), and \(\mathfrak A_\eta\neq\mathfrak A_\delta\) because
\(\mathfrak A_{\eta+1}\subseteq\mathfrak A_\delta\). Also
\(R_j^{\mathfrak A_\delta}=\bigcup_\eta R_j^{\mathfrak A_\eta}=R_j^{\mathfrak
A_0}\). Since \(\delta\) is countable, choose
\(\eta_0<\eta_1<\cdots\) cofinal in \(\delta\): enumerate
\(\delta=\{\gamma_n:n\in\omega\}\) and put \(\eta_0=\gamma_0\),
\(\eta_{n+1}=\max(\eta_n,\gamma_{n+1})+1<\delta\). Then \(\mathfrak
A_\delta=\bigcup_n\mathfrak A_{\eta_n}\) with \(\mathfrak
A_{\eta_n}\prec\mathfrak A_{\eta_{n+1}}\), and each \(\mathfrak A_{\eta_n}\),
being isomorphic to \(\mathfrak A_0\), is homogeneous and denumerable. By
Lemma~6.1(c) (\cref{thm:p03-w03-p03-029-mv-6-1}) \(\mathfrak A_\delta\) is
homogeneous and denumerable. Next, \(\mathscr I\mathscr S_\omega(\mathfrak
A_\delta^*)=\mathscr I\mathscr S_\omega(\mathfrak A_0^*)\), by
\cref{thm:p03-w03-p03-029-mv-3-1}(3): a finite non-empty
\(X\subseteq\abs{\mathfrak A_\delta}\) lies in some \(\abs{\mathfrak
A_\eta}\); with \(k\) an isomorphism of \(\mathfrak A_\eta\) onto
\(\mathfrak A_0\), \((\mathfrak A_\delta,x)_{x\in X}\equiv(\mathfrak
A_\eta,x)_{x\in X}\equiv(\mathfrak A_0,k(x))_{x\in X}\); conversely
\(\mathfrak A_0\prec\mathfrak A_\delta\), so the identity on any finite
\(X\subseteq\abs{\mathfrak A_0}\) is an elementary map from \(\mathfrak
A_0\) to \(\mathfrak A_\delta\). Finally \(\mathfrak A_\delta\equiv\mathfrak
A_0\). By Lemma~6.1(b), \(\mathfrak A_\delta\cong\mathfrak A_0\). The other
two conditions were shown above, so the recursion continues.

\emph{Step 5: the union.} Let \(\mathfrak B=\bigcup(\mathfrak
A_\xi/\xi<\omega_1)\). By the elementary chain theorem \(\mathfrak
A_0\prec\mathfrak B\), so \(\mathfrak B\equiv\mathfrak A_0\equiv\mathfrak
A_1\equiv\mathfrak A\). \(R_j^{\mathfrak B}=\bigcup_\xi R_j^{\mathfrak
A_\xi}=R_j^{\mathfrak A_0}\), so \(\overline{R_j^{\mathfrak B}}=\omega_0\).
\(\abs{\mathfrak B}\) is a union of \(\aleph_1\) countable sets, so
\(\overline{\mathfrak B}\leq\aleph_1\). Choosing for each \(\xi<\omega_1\) an
element of \(\abs{\mathfrak A_{\xi+1}}\setminus\abs{\mathfrak A_\xi}\) gives an
injection of \(\omega_1\) into \(\abs{\mathfrak B}\), since for
\(\xi<\xi'\) the element chosen for \(\xi\) lies in \(\abs{\mathfrak
A_{\xi+1}}\subseteq\abs{\mathfrak A_{\xi'}}\) and that for \(\xi'\) does not.
So \(\overline{\mathfrak B}=\omega_1\).
\end{proof}
\end{annexitem}
